% ---
% school: "college"
% level: "8eme"
% section: null
% chapter: "العمليات والحساب على الأعداد الكسرية"
% source: "Axiomaths"
% title: "الجذاء وحاصل القسمة"
% country: "Tunisie"
% year: 2026
% professor: "Axiomaths"
% difficulty: "Facile"
% points: null
% tags: ["fractions", "multiplication", "division"]
% ---
% Exercice 4/18 de la série 8ème « العمليات والحساب على الأعداد الكسرية » (statement + \axsol + solution). Compiler avec XeLaTeX + axiomath.sty (+ axiomaths-logo.png).
\documentclass[12pt,a4paper]{article}
\usepackage[level=8b,color=15803D]{axiomath}
\axsetup{title={العمليات والحساب على الأعداد الكسرية}, corrige}

\begin{document}

\begin{corr}[title={الجذاء وحاصل القسمة}, stars=2, source={Axiomaths}]
احسب ثم اختزل النتيجة إن أمكن :
\begin{enumerate}
  \item $\dfrac{14}{15}\times\dfrac{-5}{21}$
  \item $\dfrac{-8}{9}\div\dfrac{4}{27}$
  \item $\dfrac{-3}{4}\times\dfrac{-8}{15}$
  \item $\dfrac{15}{16}\div\dfrac{-5}{12}$
  \item $6\div\dfrac{-9}{4}$
  \item $\dfrac{-2}{3}\times\dfrac{-3}{2}\times\dfrac{-5}{7}$
\end{enumerate}
\axsol
\begin{prop}[الضرب]
\ib{$\dfrac{a}{b}\times\dfrac{c}{d}=\dfrac{a\times c}{b\times d}$} \qquad جذاء عددين لهما نفس الإشارة موجب ولهما إشارتان مختلفتان سالب.
\end{prop}
\begin{prop}[القسمة]
القسمة على عدد غير معدوم هي الضرب في مقلوبه : \ib{$\dfrac{a}{b}\div\dfrac{c}{d}=\dfrac{a}{b}\times\dfrac{d}{c}$}
\end{prop}
نختزل قبل الضرب لتبقى الأعداد صغيرة.
\begin{enumerate}[itemsep=14pt]
  \item الإشارة سالبة :
  \[\dfrac{14}{15}\times\dfrac{-5}{21}=-\dfrac{14}{21}\times\dfrac{5}{15}=-\dfrac{2}{3}\times\dfrac{1}{3}=-\dfrac{2}{9}\]
  \item نضرب في مقلوب $\dfrac{4}{27}$ :
  \[\dfrac{-8}{9}\div\dfrac{4}{27}=\dfrac{-8}{9}\times\dfrac{27}{4}=-\dfrac{8}{4}\times\dfrac{27}{9}=-2\times3=-6\]
  \item الإشارتان سالبتان إذن الجذاء موجب :
  \[\dfrac{-3}{4}\times\dfrac{-8}{15}=\dfrac{3}{15}\times\dfrac{8}{4}=\dfrac{1}{5}\times2=\dfrac{2}{5}\]
  \item مقلوب $\dfrac{-5}{12}$ هو $\dfrac{12}{-5}=-\dfrac{12}{5}$ :
  \[\dfrac{15}{16}\div\dfrac{-5}{12}=\dfrac{15}{16}\times\left(-\dfrac{12}{5}\right)=-\dfrac{15}{5}\times\dfrac{12}{16}=-3\times\dfrac{3}{4}=-\dfrac{9}{4}\]
  \item نكتب $6=\dfrac{6}{1}$ :
  \[6\div\dfrac{-9}{4}=6\times\dfrac{4}{-9}=-\dfrac{6}{9}\times4=-\dfrac{2}{3}\times4=-\dfrac{8}{3}\]
  \item $\dfrac{-2}{3}$ و $\dfrac{-3}{2}$ مقلوب أحدهما الآخر، إذن جذاؤهما $1$ :
  \[\dfrac{-2}{3}\times\dfrac{-3}{2}\times\dfrac{-5}{7}=1\times\dfrac{-5}{7}=-\dfrac{5}{7}\]
\end{enumerate}
\end{corr}

\end{document}
